\documentclass[final]{beamer}
\usepackage[size=custom,width=121.92,height=121.92,scale=1.75]{beamerposter}
% 121.92 cm = 48 in, the standard ISEF tri-fold display-board panel size
\usepackage[utf8]{inputenc}
\usepackage{amsmath, amssymb}
\usepackage{booktabs}
\usepackage{tikz}
\usetikzlibrary{arrows.meta, positioning, calc}

% ---------------------------------------------------------------- styling
\definecolor{myblue}{RGB}{20,60,110}
\definecolor{mylight}{RGB}{228,236,246}
\definecolor{sealcolor}{RGB}{175,45,45}
\definecolor{chaincolor}{RGB}{30,80,150}
\definecolor{freshcolor}{RGB}{40,130,70}

\setbeamercolor{block title}{bg=myblue, fg=white}
\setbeamercolor{block body}{bg=mylight, fg=black}
\setbeamertemplate{navigation symbols}{}
\setbeamertemplate{itemize items}[circle]
\setbeamerfont{block title}{size=\Large,series=\bfseries}
\setbeamerfont{block body}{size=\normalsize}
\setlength{\parskip}{6pt}

\title{An Audited Computational Framework for Knot-Invariant Inequalities}
\author{Aryan Padarthi}
\institute{Independent Research}

\begin{document}
\begin{frame}[t]
\vspace{-1.3cm}
{\centering
\color{myblue}\VERYHuge\bfseries \inserttitle \par
\vspace{10pt}
\Large\insertauthor \quad--\quad \insertinstitute \par}
\vspace{0.6cm}

\begin{columns}[t, totalwidth=\textwidth]

% ============================== COLUMN 1 ==============================
\begin{column}{.325\linewidth}

\begin{block}{Two Ways a Database Search Fails}
\textbf{Motivation.} Mining a knot-invariant database (KnotInfo, 12{,}967 knots)
for inequalities $A(K)\le B(K)$ is cheap and tempting. It fails two ways
that are easy to miss:
\begin{itemize}
\item A routine can be \textbf{structurally incapable} of finding what
  it claims to test for.
\item A bound can hold on \textbf{every tabulated example} and still
  be false beyond the database's crossing-number ceiling.
\end{itemize}
This project hit both, caught both, and built validated tools in response.
\end{block}

\begin{block}{Failure Mode 1: A Tautological Engine}
An early signature routine for positive 3-braids returned
$\sigma(\hat\beta) = -(e-2)$ \textbf{for every input}, regardless of the
braid word -- restating the crossing number $e$, never actually computing
a linking form. ``Zero defects in 1{,}280 configurations'' was forced by
the code, not discovered.

\vspace{8pt}
\textbf{Counterexample once fixed:} $T(3,7)$, a genuine positive 3-braid
knot (14 crossings, braid index 3):
\[ s = 12, \qquad |\sigma| = 8, \qquad \text{defect} = 4. \]
Re-run with a validated engine, \textbf{84--89\%} of knots in two
independent 3-braid families show a nonzero defect, at crossings as low
as 10.
\end{block}

\begin{block}{Failure Mode 2: True on Every Tabulated Knot, Still False}
The bound $g_T(K) \le \text{braid index}(K) - 1$ held with
\textbf{zero violations across all 2{,}953} KnotInfo knots with both
invariants tabulated.

\vspace{8pt}
Checked against the literature (Abe--Kishimoto 2010; Lowrance 2011)
before reporting it:
\[ g_T(T_{3,3n+i}) = n \quad (n\ge0,\ i\in\{0,1,2\}) \]
Braid index of $T(3,q)$ is fixed at 3 for $q>3$; $g_T$ grows without
bound. \textbf{First failure: $T(3,10)$}, 20 crossings, $g_T=3>2$.

\vspace{8pt}
\textbf{KnotInfo stops at 12 crossings.} The counterexample is 8
crossings past the ceiling --- an exhaustive database sweep could never
have found it.
\end{block}

\begin{block}{Also Found: Six More Silent Bugs}
Separate from the two failures above, six \texttt{print()} statements
across four analysis scripts contained hardcoded narrative claims that
did not match the data computed two lines above them in the same script
--- e.g.\ one asserted defect-knot determinants ``skew towards 1 and 7
mod 8'' when the actual dominant residues are 5 and 3 (70.8\% combined,
vs.\ 29\% for 1 and 7). All six were traced, fixed to derive their claim
from the real computation, and re-run.
\end{block}

\end{column}

% ============================== COLUMN 2 ==============================
\begin{column}{.325\linewidth}

\begin{block}{Two Validated, From-Scratch Engines}
\textbf{Seifert-matrix signature} (Collins' 2007 algorithm, positive
braids): homology generators between consecutive same-index crossings,
linking numbers by a 5-case rule. Validated against 34 torus knots
(closed-form signature) and all 17 KnotInfo knots with an authentic
positive-braid word --- \textbf{0 mismatches}.

\vspace{8pt}
\textbf{Turaev genus of a diagram} (Kauffman-state / Temperley--Lieb
circle counting, union-find implementation): validated against 3
alternating knots (must give exactly $g_T{=}0$) and \textbf{596 real
KnotInfo knots} --- never below the true tabulated value (mathematically
required), exactly tight in 25.3\% of cases.
\end{block}

\begin{block}{The Hinge Mechanism (Proof Sketch)}
\textbf{Theorem.} For a positive braid word on 3 strands with $k$
syllables (maximal same-generator blocks) and length $e$:
\[
g_T(D) = \left\lfloor \frac{k}{2} \right\rfloor - 1, \qquad k\ge2,
\]
independent of $e$ and of how crossings are split among syllables.

\vspace{14pt}
\begin{center}
\begin{tikzpicture}[scale=2.1, every node/.style={font=\normalsize}]
  \foreach \x/\lab in {0/$T_0$, 1/$T_1$, 2/$T_2$} {
    \draw[thick] (\x,0) -- (\x,4);
    \node[below] at (\x,-0.1) {\lab};
  }
  \draw[very thick, sealcolor] (0,3.3) .. controls (0.5,3.7) and (0.5,3.7) .. (1,3.3);
  \node[sealcolor, right] at (1.35,3.5) {seal at step 1};
  \draw[very thick, chaincolor, -{Stealth[length=3mm]}] (2,2.7) .. controls (1.4,2.3) .. (1,1.9);
  \draw[very thick, chaincolor, -{Stealth[length=3mm]}] (1,1.9) .. controls (1.4,1.5) .. (2,1.1);
  \node[chaincolor, right] at (2.35,2.3) {chain absorbs};
  \node[chaincolor, right] at (2.35,1.5) {via shared hinge};
  \filldraw[freshcolor] (2,0.5) circle (5pt);
  \node[freshcolor, right] at (2.35,0.5) {latest fresh node};
\end{tikzpicture}
\end{center}

\textbf{Lemma 1 (reduction).} Repeats within a syllable spin off isolated
circles; only one crossing per syllable matters. Reduces to $k$-length
pure alternation.

\vspace{6pt}
\textbf{Lemma 2 (hinge induction).} $\sigma_1{=}\mathrm{cap}(0,1)$ and
$\sigma_2{=}\mathrm{cap}(1,2)$ share strand 1. Before closure: exactly 3
components always --- $\{T_0,T_1\}$ sealed at step 1, a chain absorbing
$T_2$ one node per step, and the newest fresh node. Closure parity gives
2 components ($k$ odd) or 1 ($k$ even).

\vspace{6pt}
Zero mismatches: base case $k{=}1..40$, 160 random distributions, 203
further $(e,k)$ pairs.
\end{block}

\begin{block}{Why 8$_{19}$ Motivated This}
KnotInfo's own word for $8_{19}=T(3,4)$, \texttt{[1,1,1,2,1,1,1,2]}
($k{=}4$ syllables), gives $g_T(D)=1$ --- the true minimum. The naive
torus word \texttt{[1,2,1,2,1,2,1,2]} ($k{=}8$) for the \emph{same knot,
same crossing number, same braid index} gives $g_T(D)=3$. Syllable count,
not crossing number, controls diagram genus.
\end{block}

\end{column}

% ============================== COLUMN 3 ==============================
\begin{column}{.325\linewidth}

\begin{block}{Correction: Not Original --- Caught the Same Way}
\centering
{\color{sealcolor}\Large\textbf{Checked against literature before reporting. It failed.}}

\vspace{8pt}
\raggedright
A. Lowrance (\emph{Comment. Math. Helv.} 86, 2011) already proves the
reduction lemma in general (\textbf{Corollary 3.10}, any braid index),
building on Champanerkar--Kofman. His Proposition 4.3 even uses the
even-$k$ base case as an intermediate step.

\vspace{6pt}
\textbf{What survives:} an independent proof by a different technique
(union-find/Temperley-Lieb vs.\ his tangle-twisting argument), the
explicit closed form, and the validated reusable engine.

\vspace{6pt}
This is the \textbf{third} time in this project that checking a claim
against data or literature before reporting it changed the conclusion.
\end{block}

\begin{block}{Beyond 3 Strands: Labeled a Conjecture, Not a Theorem}
Cyclic alternation for $n=3..9$ strands:
\[
\text{period} = \begin{cases}2 & n{-}1 \text{ even}\\ n{-}1 & n{-}1\text{ odd}\end{cases}
\]
\[
\text{growth per syllable} = \frac{\lfloor (n-1)/2\rfloor}{n-1}
\]
Consistent with 2-colorability of the generator cycle $C_{n-1}$ --- not
proved. \textbf{No literature check done yet; no proof attempted for
$n\ge5$.} Reported honestly as future work, not a result.
\end{block}

\begin{block}{Results at a Glance}
\begin{center}
\small
\begin{tabular}{@{}p{0.40\linewidth}p{0.20\linewidth}p{0.32\linewidth}@{}}
\toprule
\textbf{Claim} & \textbf{Status} & \textbf{Evidence} \\
\midrule
Positive 3-braids signature-rigid & \textbf{Retracted} & tautological engine \\[4pt]
$g_T\le$ braid index $- 1$ & \textbf{Refuted} & $T(3,10)$ + cited thm. \\[4pt]
Syllable formula ($n{=}3$) & Re-proved & Lowrance 2011 (prior) \\[4pt]
Syllable formula ($n{\ge}4$) & Conjecture & computation only \\
\bottomrule
\end{tabular}
\end{center}
\end{block}

\begin{block}{Discussion \& Future Work}
The most reliable output of this project was not a single inequality but
a \textbf{discipline}: re-run from a clean checkout and diff against the
claim; check any surviving claim against the literature before reporting
it. Applied consistently, it changed the conclusion three separate times.

\vspace{6pt}
\textbf{Next:} literature search for the $n\ge4$ conjecture; a proof via
the hinge structure of the generator cycle $C_{n-1}$; extend both engines
to mixed-sign (non-positive) braid words.

\vspace{10pt}
\small
Full reproducible record: \texttt{manuscript/log.tex} (20 dated entries),
every number traceable to a script and a saved file under
\texttt{results/}. Paper: \texttt{manuscript/isef\_paper.tex}.
\end{block}

\end{column}

\end{columns}
\end{frame}
\end{document}
